\documentclass{fiche}
\metadonnees{13}{La quatrième dimension}{Bloc 3 — La ville, le travail, la science}

\begin{document}
\entete

\objectifs{%
\textbf{Langue} : propositions relatives, déterminatives et appositives ; relative sans
pronom.\par
\textbf{Texte} : H.\,G. Wells, \emph{The Time Machine} (1895), chapitre \textsc{i}.\par
\textbf{Durée} : 1\,h\,30 — exercices 35 min, texte et questions 30 min, version 25 min.}

%% ===================================================================
\exercice[13 min]{Le pronom relatif}

\rappel{\textbf{Rappel.} \emph{Who} pour les personnes, \emph{which} pour les choses,
\emph{that} pour les deux. Quand le relatif est complément du verbe, il peut disparaître.}

\consigne{Complétez avec \emph{who}, \emph{which}, \emph{that}, \emph{whose}, ou laissez vide
(Ø) si le pronom peut s'effacer.}

\begin{anglais}
\begin{enumerate}[label=\arabic*.,itemsep=2.2mm,leftmargin=*]
  \item The geometry \trou{Ø} they taught you at school is founded on a misconception.
  \item A line \trou{which / that} has no thickness has no real existence.
  \item Filby, \trou{who} had red hair, was an argumentative person.
  \item The chairs, \trou{which} were his own patents, embraced us rather than submitted.
  \item Some people \trou{who / that} talk about the Fourth Dimension do not know what they
    mean.
  \item This is the diagram \trou{Ø} I traced with my finger.
  \item The man \trou{whose} portrait you see was eight years old.
  \item The direction \trou{in which} our consciousness moves is called Time.
  \item Everything \trou{that} he said seemed absurd and was true.
  \item The Provincial Mayor, \trou{whose} brows were knitted, said nothing.
  \item There is no difference \trou{Ø} we can measure between Time and Space.
  \item The dimension \trou{Ø} they had forgotten was Duration.
  \item Newcomb, \trou{who} was expounding this in New York, is a real mathematician.
  \item Any body \trou{which / that} is real must have extension in four directions.
\end{enumerate}
\end{anglais}

\note{Le pronom s'efface en 1, 6, 11, 12 : il y est complément, et un sujet le suit
immédiatement. Il ne peut jamais s'effacer quand il est sujet (2, 5, 9, 14) ni dans une
relative entre virgules (3, 4, 10, 13).}

\note[Le test de l'effacement]{Après le nom, s'il vient directement un autre groupe nominal
ou un pronom (\emph{the geometry they taught}), c'est que le relatif était complément :
il peut disparaître. S'il vient un verbe conjugué (\emph{a line has no thickness}), il était
sujet : il reste obligatoire. C'est la faute de lecture la plus fréquente en version ---
on cherche le verbe de la principale et on trouve celui de la relative.}

%% ===================================================================
\exercice[10 min]{Déterminative ou appositive}

\rappel{\textbf{Rappel.} La relative déterminative identifie le nom et ne prend pas de
virgule ; l'appositive ajoute un renseignement et s'encadre de virgules. \emph{That} est
impossible dans une appositive.}

\consigne{Ajoutez les virgules quand il le faut, et dites ce que la relative change au sens.}

\begin{anglais}
\begin{enumerate}[label=\arabic*.,itemsep=2.5mm,leftmargin=*]
  \item The guests who had dined well listened politely.
    \lignes{2}
    \corr{Sans virgules : seuls ceux qui avaient bien dîné écoutaient. Avec virgules : tous
    avaient bien dîné, et tous écoutaient.}
  \item Filby who had red hair objected at once.
    \lignes{2}
    \corr{Filby, who had red hair, objected at once. Appositive : il n'y a qu'un Filby.}
  \item A cube which does not last for any time cannot exist.
    \lignes{2}
    \corr{Pas de virgules : la relative définit de quel cube on parle.}
  \item The Time Traveller whose chairs were patented smiled.
    \lignes{2}
    \corr{The Time Traveller, whose chairs were patented, smiled. Appositive.}
  \item Scientific people who know this very well say nothing.
    \lignes{2}
    \corr{Ambigu : sans virgules, seuls certains savants ; avec virgules, tous les savants.}
  \item Mathematics which he had studied at Cambridge bored him.
    \lignes{2}
    \corr{Mathematics, which he had studied at Cambridge, bored him. Une seule mathématique :
    appositive.}
\end{enumerate}
\end{anglais}

\note[La virgule change le sens]{\emph{My brother who lives in Leeds is a doctor} sous-entend
que j'ai plusieurs frères ; \emph{My brother, who lives in Leeds, is a doctor} sous-entend que
je n'en ai qu'un. Le français, lui, marque la différence à l'oral par l'intonation et à
l'écrit par les mêmes virgules : la règle se transpose sans perte.}

%% ===================================================================
\exercice[12 min]{Le vocabulaire de la démonstration}
\consigne{a) Associez chaque mot à sa traduction.}

\begin{multicols}{2}
\begin{enumerate}[label=\arabic*.,itemsep=1.2mm,leftmargin=*]
  \item to expound \hfill \trou{d}
  \item recondite \hfill \trou{j}
  \item a misconception \hfill \trou{b}
  \item thickness \hfill \trou{g}
  \item breadth \hfill \trou{a}
  \item a plane \hfill \trou{l}
  \item to overlook \hfill \trou{c}
  \item to controvert \hfill \trou{i}
  \item a ground \hfill \trou{f}
  \item to admit \hfill \trou{k}
  \item a statement \hfill \trou{e}
  \item at right angles \hfill \trou{h}
\end{enumerate}
\columnbreak
\begin{enumerate}[label=\alph*.,itemsep=1.2mm,leftmargin=*]
  \item la largeur
  \item une idée fausse
  \item négliger, ne pas voir
  \item exposer, développer
  \item un énoncé, une affirmation
  \item un fondement, une raison
  \item l'épaisseur
  \item perpendiculairement
  \item réfuter, contester
  \item abscons, ardu
  \item concéder, reconnaître
  \item un plan (géométrie)
\end{enumerate}
\end{multicols}

\note[Trois dimensions, trois noms]{\emph{long} $\rightarrow$ \emph{length},
\emph{broad} $\rightarrow$ \emph{breadth}, \emph{thick} $\rightarrow$ \emph{thickness}. Les
deux premiers changent de voyelle, comme \emph{strong} $\rightarrow$ \emph{strength},
\emph{wide} $\rightarrow$ \emph{width}, \emph{deep} $\rightarrow$ \emph{depth}.}

\consigne{b) Complétez avec un mot de la liste.}
\begin{anglais}
\begin{enumerate}[label=\arabic*.,itemsep=2mm,leftmargin=*]
  \item A mathematical line has no \trou{thickness} and no real existence.
  \item He was \trou{expounding} a \trou{recondite} matter to his guests.
  \item The geometry they taught you is founded on a \trou{misconception}.
  \item Space is definable by three planes, each \trou{at right angles} to the others.
\end{enumerate}
\end{anglais}

%% ===================================================================
\newpage
\section{Texte}

\noindent\small\textit{Un soir, après dîner, chez un inventeur que le narrateur appelle
seulement \og le Voyageur du Temps \fg.}\normalsize

\begin{texte}
The Time Traveller (for so it will be convenient to speak of him) was
\gl[to expound]{expounding}{exposer, développer} a \gl[recondite]{recondite}{abscons, ardu}
matter to us. His grey eyes shone and twinkled, and his usually pale face was flushed and
animated. The fire burnt brightly, and the soft radiance of the incandescent lights in the
lilies of silver caught the bubbles that flashed and passed in our glasses. Our chairs, being
his patents, embraced and caressed us rather than submitted to be sat upon, and there was
that luxurious after-dinner atmosphere, when thought runs gracefully free of the
\gl[trammels]{trammels}{les entraves} of precision. And he put it to us in this way ---
marking the points with a lean \gl[a forefinger]{forefinger}{l'index} --- as we sat and lazily
admired his \gl[earnestness]{earnestness}{le sérieux, la gravité} over this new paradox (as we
thought it) and his \gl[fecundity]{fecundity}{la fécondité}.

``You must follow me carefully. I shall have to \gl[to controvert]{controvert}{réfuter,
contester} one or two ideas that are almost universally accepted. The geometry, for instance,
they taught you at school is founded on a \gl[a misconception]{misconception}{une idée
fausse}.''

``Is not that rather a large thing to expect us to begin upon?'' said Filby, an
argumentative person with red hair.

``I do not mean to ask you to accept anything without reasonable ground for it. You will soon
admit as much as I need from you. You know of course that a mathematical line, a line of
thickness \emph{nil}, has no real existence. They taught you that? Neither has a mathematical
plane. These things are mere abstractions.''

``That is all right,'' said the Psychologist.

``Nor, having only length, \gl[breadth]{breadth}{la largeur} and thickness, can a cube have a
real existence.''

``There I object,'' said Filby. ``Of course a solid body may exist. All real things ---''

``So most people think. But wait a moment. Can an \emph{instantaneous} cube exist?''

``Don't follow you,'' said Filby.

``Can a cube that does not last for any time at all, have a real existence?''

Filby became pensive. ``Clearly,'' the Time Traveller proceeded, ``any real body must have
extension in \emph{four} directions: it must have Length, Breadth, Thickness, and ---
Duration. But through a natural \gl[an infirmity]{infirmity}{une faiblesse, une infirmité} of
the flesh, which I will explain to you in a moment, we incline to
\gl[to overlook]{overlook}{négliger, ne pas voir} this fact. There are really four
dimensions, three which we call the three planes of Space, and a fourth, Time. There is,
however, a tendency to draw an unreal distinction between the former three dimensions and the
latter, because it happens that our consciousness moves intermittently in one direction along
the latter from the beginning to the end of our lives.''

``That,'' said a very young man, making \gl[spasmodic]{spasmodic}{saccadé} efforts to relight
his cigar over the lamp; ``that\ldots{} very clear indeed.''

``Now, it is very remarkable that this is so extensively overlooked,'' continued the Time
Traveller, with a slight \gl[an accession]{accession}{un surcroît} of cheerfulness. ``Really
this is what is meant by the Fourth Dimension, though some people who talk about the Fourth
Dimension do not know they mean it. It is only another way of looking at Time.
\emph{There is no difference between Time and any of the three dimensions of Space except
that our consciousness moves along it}. But some foolish people have got hold of the wrong
side of that idea.''
\end{texte}
\source{H.\,G. Wells, \emph{The Time Machine}, Londres, William Heinemann, 1895,
ch.~\textsc{i}.}

\partiel[15 min]{Questions}
\consigne{\en{Answer in English, in full sentences.}}

\begin{enumerate}[label=\arabic*.,itemsep=2mm,leftmargin=*]
  \item \en{Describe the setting of the scene. Why does it matter?}
    \lignes{3}
    \corr{A comfortable house after dinner: the fire burning, wine in the glasses, patented
    chairs that embrace the guests. It matters because it is the moment when \emph{thought
    runs gracefully free of the trammels of precision} --- the guests are too comfortable to
    argue properly, which is exactly what the Time Traveller needs.}
  \item \en{What does the Time Traveller ask his guests to accept first?}
    \lignes{2}
    \corr{That a mathematical line, having no thickness, has no real existence, and neither
    has a mathematical plane: they are mere abstractions.}
  \item \en{What is his argument about the cube?}
    \lignes{3}
    \corr{If a line without thickness cannot exist, then a cube with only length, breadth and
    thickness cannot either, since an instantaneous cube --- one that lasts no time at all
    --- is impossible. A real body must therefore also have duration.}
  \item \en{According to him, why do we treat Time as different from the other three
    dimensions?}
    \lignes{3}
    \corr{Because of \emph{a natural infirmity of the flesh}: our consciousness moves along
    Time in one direction only, from birth to death, so we mistake our own limitation for a
    property of Time.}
  \item \en{How do the guests react? Look closely at the young man and at Filby.}
    \lignes{3}
    \corr{Filby objects twice but is answered each time and gives up. The young man says
    \emph{that\ldots{} very clear indeed} while struggling to relight his cigar, which shows
    that he has understood nothing and is pretending. Nobody really resists.}
  \item \en{Where do you find the narrator's irony?}
    \lignes{2}
    \corr{In the parenthesis \emph{as we thought it}, which admits that they were wrong; in
    \emph{lazily admired}; in the chairs that \emph{embraced and caressed} their occupants.
    The narrator is judging his own comfortable scepticism after the event.}
\end{enumerate}

\note[Les relatives du texte]{\emph{the bubbles that flashed and passed} (déterminative),
\emph{Our chairs, being his patents,} (participe apposé, équivalent d'une appositive),
\emph{ideas that are almost universally accepted}, \emph{a cube that does not last},
\emph{three which we call the three planes of Space}, et surtout \emph{The geometry\ldots
they taught you at school} : relative sans pronom, avec en plus un incise entre le nom et la
relative. C'est la phrase la plus difficile à lire du texte.}

%% ===================================================================
\newpage
\section{Version}
\consigne{Traduisez le passage en français. Il suit immédiatement le texte.}

\begin{version}
``Well, I do not mind telling you I have been at work upon this geometry of Four Dimensions
for some time. Some of my results are curious. For instance, here is a portrait of a man at
eight years old, another at fifteen, another at seventeen, another at twenty-three, and so
on. All these are evidently sections, as it were, Three-Dimensional representations of his
Four-Dimensioned being, which is a fixed and \gl[unalterable]{unalterable}{immuable} thing.''

``Scientific people,'' proceeded the Time Traveller, after the pause required for the proper
\gl[assimilation]{assimilation}{l'assimilation} of this, ``know very well that Time is only a
kind of Space. Here is a popular scientific diagram, a weather record. This line I trace with
my finger shows the movement of the barometer. Yesterday it was so high, yesterday night it
fell, then this morning it rose again, and so gently upward to here. Surely the
\gl[mercury]{mercury}{le mercure} did not trace this line in any of the dimensions of Space
generally recognised? But certainly it traced such a line, and that line, therefore, we must
conclude, was along the Time-Dimension.''

``But,'' said the Medical Man, staring hard at a coal in the fire, ``if Time is really only a
fourth dimension of Space, why is it, and why has it always been, regarded as something
different? And why cannot we move about in Time as we move about in the other dimensions of
Space?''
\end{version}
\source{\emph{Ibid.}}

\lignes{22}

\corr{\textbf{Proposition de traduction.} \og Eh bien, je ne vois pas d'inconvénient à vous
dire que je travaille depuis quelque temps à cette géométrie des quatre dimensions. Certains
de mes résultats sont curieux. Tenez, voici le portrait d'un homme à huit ans, un autre à
quinze, un autre à dix-sept, un autre à vingt-trois, et ainsi de suite. Tous sont évidemment
des coupes, pour ainsi dire, des représentations à trois dimensions de son être à quatre
dimensions, lequel est une chose fixe et immuable. \fg

\og Les hommes de science, reprit le Voyageur du Temps après la pause qu'exigeait la bonne
assimilation de ceci, savent fort bien que le Temps n'est qu'une sorte d'espace. Voici un
diagramme scientifique bien connu, un relevé météorologique. Cette ligne, que je suis du
doigt, montre le mouvement du baromètre. Hier il était à telle hauteur, hier soir il est
descendu, puis ce matin il est remonté, et doucement jusqu'ici. Assurément le mercure n'a pas
tracé cette ligne dans aucune des dimensions de l'espace communément reconnues ? Et pourtant
il a bien tracé une telle ligne, et cette ligne, il nous faut donc en conclure, était le long
de la dimension du Temps. \fg

\og Mais, dit le Médecin, les yeux fixés sur un charbon du feu, si le Temps n'est réellement
qu'une quatrième dimension de l'espace, pourquoi est-il, et pourquoi a-t-il toujours été,
considéré comme quelque chose de différent ? Et pourquoi ne pouvons-nous pas nous déplacer
dans le Temps comme nous nous déplaçons dans les autres dimensions de l'espace ? \fg}

\note[Choix de traduction 1]{\emph{I have been at work upon this geometry for some time} :
present perfect en \emph{be + -ing} avec \emph{for}, donc présent français + \og depuis \fg.
La faute de version classique est le passé composé.}

\note[Choix de traduction 2]{\emph{This line I trace with my finger} : relative sans pronom,
objet antéposé. Le français doit rétablir le relatif (\og que je suis du doigt \fg). Même
structure qu'à l'exercice 1.}

\note[Choix de traduction 3]{\emph{which is a fixed and unalterable thing} : appositive
portant sur tout le groupe précédent. \og Lequel \fg{} lève l'ambiguïté que \og qui \fg{}
laisserait planer sur l'antécédent.}

\note[Choix de traduction 4]{\emph{after the pause required for the proper assimilation of
this} : participe passé épithète, très concis. Le français déplie en relative (\og la pause
qu'exigeait\ldots \fg) plutôt que de calquer \og requise pour \fg, qui sonne administratif.}

\note[Choix de traduction 5]{\emph{why is it, and why has it always been, regarded as} : un
seul participe pour deux auxiliaires. Le français doit répéter le verbe : \og pourquoi
est-il, et pourquoi a-t-il toujours été, considéré \fg.}

%% ===================================================================
\partiel{Lexique à mémoriser}
\begin{itemize}[label=--,itemsep=0.8mm,leftmargin=*]\small
  \item \textbf{length} — la longueur ; long
  \item \textbf{breadth} — la largeur ; broad
  \item \textbf{width} — la largeur ; wide
  \item \textbf{thickness} — l'épaisseur ; thick
  \item \textbf{depth} — la profondeur ; deep
  \item \textbf{strength} — la force ; strong
  \item \textbf{a plane} — un plan ; \textit{et} un avion
  \item \textbf{a statement} — un énoncé ; \textit{verbe} to state \textit{(fiche 11)}
  \item \textbf{a ground} — un fondement ; on what grounds? \textit{« à quel titre ? »}
  \item \textbf{to overlook} — négliger ; \textit{et} donner sur \textit{(une vue)}
  \item \textbf{to admit} — reconnaître, concéder ; \textit{et} admettre
  \item \textbf{evidently} — manifestement ; \textit{et non « évidemment » au sens de « bien sûr »}
  \item \textbf{eventually} — finalement \textit{(fiche 5)}
  \item \textbf{mere} — simple, pur ; merely \textit{« simplement »}
  \item \textbf{to last} — durer ; \textit{et} dernier
  \item \textbf{a record} — un relevé, un registre ; \textit{verbe} to record
  \item \textbf{to trace} — tracer, suivre à la trace
  \item \textbf{to regard as} — considérer comme ; \textit{à ne pas confondre avec} to look at
  \item \textbf{lean} — maigre, sec ; \textit{verbe} to lean \textit{« s'appuyer »}
  \item \textbf{flushed} — rouge, empourpré ; \textit{de} to flush \textit{« rougir »}
\end{itemize}

\end{document}
